% =====================================================================
%  Archimedean first-slab positivity for shifted zeta canonical systems
%  An off-center Weil generator and a radial energy identity
%  Preprint v1.0 -- 5 September 2026 -- Edward Baker
%  License: CC BY 4.0
% =====================================================================
\documentclass[11pt,a4paper]{article}

\usepackage[utf8]{inputenc}
\usepackage[protrusion=true,expansion=false]{microtype}
\usepackage[margin=28mm]{geometry}
\usepackage{amsmath,amssymb,amsthm,mathtools}
\usepackage{booktabs}
\usepackage{array}
\usepackage{enumitem}
\setlist[itemize]{label=$\bullet$}
\usepackage[hidelinks]{hyperref}
\usepackage{url}
\makeatletter\g@addto@macro{\UrlBreaks}{\do\_}\makeatother
\usepackage{xcolor}

\hypersetup{
  pdftitle={Archimedean first-slab positivity for shifted zeta canonical systems},
  pdfauthor={Edward Baker},
  pdfsubject={Computer-assisted positivity certificate for the prime-free first slab of Suzuki's shifted zeta Hankel system},
  pdfkeywords={Riemann zeta function, Weil positivity, canonical systems, Hankel operators, Volterra operators, computer-assisted proof, interval arithmetic}
}

% ---------------------------------------------------------------- macros
\newcommand{\R}{\mathbb R}
\newcommand{\C}{\mathbb C}
\newcommand{\Z}{\mathbb Z}
\newcommand{\Cp}{\mathbb C_+}
\newcommand{\Ginf}{\Gamma_{\!\infty}}
\newcommand{\Dlog}{\mathcal D_{\log}}
\renewcommand{\Re}{\operatorname{Re}}
\renewcommand{\Im}{\operatorname{Im}}
\newcommand{\supp}{\operatorname{supp}}
\newcommand{\dd}{\,\mathrm d}
\newcommand{\eps}{\varepsilon}
\newcommand{\Cc}{\mathcal C}
\newcommand{\Ss}{\mathcal S}
\newcommand{\Hcal}{\mathcal H}
\newcommand{\Gcal}{\mathcal G}
\newcommand{\Kfrak}{\mathfrak K}
\newcommand{\Hsf}{\mathsf H}
\newcommand{\LZ}{L_{\mathrm Z}}
\newcommand{\pt}{\mathrm{pt}}
\DeclarePairedDelimiter{\abs}{\lvert}{\rvert}
\DeclarePairedDelimiter{\norm}{\lVert}{\rVert}
\newcommand{\succposdef}{\succ0}

% ---------------------------------------------------------------- theorems
\theoremstyle{plain}
\newtheorem{theorem}{Theorem}[section]
\newtheorem{proposition}[theorem]{Proposition}
\newtheorem{lemma}[theorem]{Lemma}
\newtheorem{corollary}[theorem]{Corollary}
\newtheorem*{thmA}{Theorem A}
\newtheorem*{thmB}{Theorem B}
\newtheorem*{thmC}{Theorem C}
\theoremstyle{definition}
\newtheorem{definition}[theorem]{Definition}
\newtheorem{remark}[theorem]{Remark}

\numberwithin{equation}{section}

\title{\textbf{Archimedean first-slab positivity for shifted zeta canonical systems}\\[1ex]
\large An off-center Weil generator and a radial energy identity}
\author{Edward Baker\\[0.5ex]
\normalsize \texttt{edwardbaker86@gmail.com}}
\date{Preprint, version 1.0 --- 5 September 2026}

\begin{document}
\maketitle

\begin{center}
\begin{minipage}{0.9\textwidth}
\small\itshape
This manuscript reports computer-assisted mathematics and was prepared with
substantial language-model assistance. The author is responsible for its
content. The analytic inputs, the interval-arithmetic certificate, and the
independent audit are described in Sections~7 and~8 and Appendix~B; the
present verification status of every component, including the items that
have not yet received line-by-line human verification, is recorded in
Section~13 and Appendix~C. The preprint is deposited to fix a citable
version of the argument while that verification proceeds.
\end{minipage}
\end{center}

\begin{abstract}
For $\omega>0$, Suzuki's shifted scattering function
$\Theta_\omega(z)=\xi(\tfrac12-\omega-iz)/\xi(\tfrac12+\omega-iz)$ admits a
prime--gamma Hankel kernel. By results of Suzuki, $\Theta_\omega$ is a
meromorphic inner function in the upper half-plane exactly when the associated
Hankel transform is an isometry, and innerness for every $\omega>\omega_0$ is
equivalent to the zero-free half-plane $\Re s>\tfrac12+\omega_0$; the Riemann
hypothesis is the case of all $\omega_0>0$. Contractivity of every finite
section of the Hankel operator is therefore a necessary condition for these
open assertions when $0<\omega<\tfrac12$.

We isolate the first logarithmic radial slab, before the first prime layer
appears. On a slab of length $L\le\log2$ the arithmetic Hankel operator reduces
exactly to a reflected finite-horizon Volterra convolution
$H_{\omega,L}=V_{\omega,L}R_L$ with completed-gamma transfer
$K_\omega(p)=\Ginf(\tfrac12+p-\omega)/\Ginf(\tfrac12+p+\omega)$.
Differentiating this transfer in the shift produces a self-adjoint generator
$Q_{\omega,L}$, an off-center prime-free Weil form, with the exact
decomposition $Q_{\omega,L}=Q_{0,L}+C_{\omega,L}$, where $Q_{0,L}$ is the
central compact-window Weil form and $C_{\omega,L}$ is the compact integral
operator with kernel
$c_\omega(x,y)=(\cosh(\omega\abs{x-y})-1)\bigl(e^{\abs{x-y}/2}-e^{-5\abs{x-y}/2}/(1-e^{-2\abs{x-y}})\bigr)$.
The apparent singularity at $\omega=\tfrac12$ cancels in this representation
and all generators share one logarithmic form domain. An Arb-certified
Legendre calculation, closed by analytic Taylor and infinite-tail bounds,
proves $Q_{\omega,L}\succ0$ for $0\le\omega\le\tfrac12$ and $0<L\le\log2$,
with an explicit coercivity constant $6.1159\times10^{-4}$ on the whole form
domain. The certificate uses $256$ shift boxes in each parity sector and
sixteen Legendre modes per sector; a clean-room reimplementation reproduces
all $512$ boxes.

We then prove a strong-operator energy identity expressing the defect
$I-V_{\omega,L}^{*}V_{\omega,L}$ as $2\int_0^\omega V_{s,L}^{*}Q_{s,L}V_{s,L}\dd s$.
Strict generator positivity, the strong limit $V_{s,L}\to I$ as $s\downarrow0$,
and compactness of $V_{\omega,L}$ give $\norm{V_{\omega,L}}<1$ and
$I\pm H_{\omega,L}\succ0$ throughout the first slab for $0<\omega\le\tfrac12$.
This complements Suzuki's unconditional strict contractivity of all finite
sections for $\omega>\tfrac12$ and settles the critical shift interval on the
prime-free slab. The energy identity also exhibits an explicit shift-radial
Gram factor with local density $Q_{s,L}$. The result is prime-free and
finite-depth: it proves neither the Riemann hypothesis nor a new zero-free
region. Its content is an independently positive archimedean base slab for the
shifted-zeta canonical-system program.
\end{abstract}

\medskip
\noindent\textbf{MSC 2020.} 11M26, 47B35, 47G10, 65G20, 65G30.\\
\textbf{Keywords.} Riemann zeta function; Weil positivity; canonical systems;
Hankel and Volterra operators; reflection positivity; computer-assisted proof;
interval arithmetic.

\setcounter{tocdepth}{1}
\tableofcontents

% =====================================================================
\section{Introduction}
% =====================================================================

\subsection{Shifted zeta positivity}

Let
\begin{equation}
\xi(s)=\tfrac12 s(s-1)\pi^{-s/2}\Gamma(s/2)\zeta(s)
\end{equation}
be the completed Riemann xi-function and, following Suzuki~\cite{Suzuki2012},
put for $\omega>0$
\begin{equation}\label{eq:Theta}
\Theta_\omega(z)=\frac{\xi(\tfrac12-\omega-iz)}{\xi(\tfrac12+\omega-iz)}.
\end{equation}
The functional equation $\xi(s)=\xi(1-s)$ gives, for real $u$,
\begin{equation}
\Theta_\omega(u)=\frac{\xi(\tfrac12+\omega+iu)}{\xi(\tfrac12+\omega-iu)},
\qquad \abs{\Theta_\omega(u)}=1 .
\end{equation}
This boundary unitarity is unconditional. The nontrivial condition is
\emph{meromorphic innerness} in the upper half-plane
$\Cp=\{\Im z>0\}$: in the terminology of~\cite{Suzuki2012}, a bounded
analytic function on $\Cp$ with unimodular boundary values almost everywhere
which extends meromorphically to $\C$.

We pin the relation to zero-free regions to the statements actually proved
in~\cite{Suzuki2012}, since it carries the significance claim of this paper.
Suzuki proves~\cite[Prop.~1.2]{Suzuki2012} that for $\omega_0>0$ the following
are equivalent: (i) $\zeta(s)\neq0$ for $\Re s>\tfrac12+\omega_0$;
(ii) $\Theta_\omega$ is a meromorphic inner function in $\Cp$ for every
$\omega>\omega_0$. In particular the Riemann hypothesis is equivalent to
meromorphic innerness of $\Theta_\omega$ for every $\omega>0$, and for
$0<\omega_0<\tfrac12$ the half-plane statement (i) is an open quasi-RH
assertion. The equivalence is stated for the whole family
$\omega>\omega_0$; for a single shift, innerness of $\Theta_\omega$ excludes
every zero $\rho$ of $\xi$ with $\Re\rho>\tfrac12+\omega$ whose pole
$z=i(\rho-\tfrac12-\omega)$ of~\eqref{eq:Theta} is not cancelled by a zero of
the numerator, and the passage to all $\omega>\omega_0$ removes the
cancellation cases. Suzuki further shows~\cite[Thm.~2.2]{Suzuki2012} that
$\Theta_\omega$ is meromorphic inner if and only if the Hankel operator
$\Hsf_\omega$ of Section~2 maps $L^2((1,\infty),dx)$ into
$L^2((0,\infty),dx)$---equivalently, if and only if on $L^2((1,\infty),dx)$
the multiplier action of $\Theta_\omega$ coincides with $\Hsf_\omega$---and
that in this case $\Hsf_\omega$ is an isometry of $L^2((0,\infty),dx)$
\cite[Lemma~4.1]{Suzuki2012}. Every compression of an isometry is a
contraction. Hence contractivity of
every finite section of $\Hsf_\omega$ is a necessary condition for innerness,
and thus, for $0<\omega<\tfrac12$, for the open zero-free statements; compare
the classical positivity formulations of Weil and Lagarias~\cite{Lagarias1999}.
Theorem~A verifies this necessary condition unconditionally on the entire
first slab and throughout the critical shift interval $0<\omega\le\tfrac12$.
For $\omega>\tfrac12$ the zero-free statement follows from the classical
critical strip, and Suzuki~\cite[Lemma~4.4]{Suzuki2012} proves the strict
finite-section contractivity there; see Section~9.3. The genuinely
RH-sensitive shift range is the one treated here.

The global problem is arithmetic. Suzuki's Hankel kernel contains translated
copies of one archimedean impulse at logarithmic depths $\log n$, weighted by
Euler-product coefficients. Proving compatible contractivity at every depth
for shifts accumulating at zero is an RH-level problem.

The present paper addresses a local but independently specified part of that
system. Before the first threshold $\log2$, no prime layer is present. The
kernel is determined solely by the completed gamma factor, so its positivity
can be studied without zeta zeros, a zero-free hypothesis, or an assumed
positive spectral measure.

\subsection{The first slab and the proof architecture}

An earlier approach to the first slab proceeded through the singular limit
$\omega\downarrow0$, compact Mosco convergence to a prime-free Weil form, and
direct parameter enclosures for the defect channels. That route remains a
useful cross-check (Section~11), but it obscures a simpler monotonicity
mechanism.

The key observation is to differentiate the causal transfer in $\omega$. If
\begin{equation}
\partial_\omega V_{\omega,L}=-A_{\omega,L}V_{\omega,L},
\end{equation}
then the defect $D_{\omega,L}=I-V_{\omega,L}^{*}V_{\omega,L}$ satisfies
\begin{equation}
\partial_\omega D_{\omega,L}
=V_{\omega,L}^{*}\bigl(A_{\omega,L}+A_{\omega,L}^{*}\bigr)V_{\omega,L}.
\end{equation}
Thus contractivity follows if the self-adjoint shift generator is positive.
The generator is not a new unrelated form: it is an off-center deformation of
the same compact-window Weil form that appears at $\omega=0$. More precisely,
after the normalization
\begin{equation}\label{eq:Qdef-intro}
Q_{\omega,L}:=\tfrac12\bigl(A_{\omega,L}+A_{\omega,L}^{*}\bigr),
\end{equation}
all shift dependence is a small bounded compact kernel. This is the structural
reason the full shift interval is certifiable with a modest matrix
calculation.

\subsection{Main results}

Let $R_Lf(x)=f(L-x)$ on $L^2(0,L)$. Let $V_{\omega,L}$ be the causal Volterra
operator defined in Section~2, and put $H_{\omega,L}=V_{\omega,L}R_L$.

\begin{thmA}[critical-shift first-slab positivity]
For every
\[
0<\omega\le\tfrac12,\qquad 0<L\le\log2,
\]
one has
\begin{equation}
\norm{V_{\omega,L}}=\norm{H_{\omega,L}}<1,
\qquad\text{and therefore}\qquad
I\pm H_{\omega,L}\succ0 .
\end{equation}
\end{thmA}

Theorem~A follows from Theorems~B and~C below. Section~9.3 records the
relation to Suzuki's complementary range $\omega>\tfrac12$ without using that
range in Theorem~A. The computer-assisted input is isolated in one theorem.

\begin{thmB}[computer-assisted off-center Weil positivity]
Let $0\le\omega\le\tfrac12$ and $0<L\le\log2$. Then the closed form
$Q_{\omega,L}$ with logarithmic form domain $\Dlog(I_L)$ (Section~4.4) is
strictly positive. More precisely, for every $f\in\Dlog(I_L)$,
\begin{equation}\label{eq:thmB-coercive}
Q_{\omega,L}[f]\ \ge\ c_*\norm f_2^2,
\qquad c_*=6.1159\times10^{-4},
\end{equation}
and $c_*$ may be replaced by $5.2613\times10^{-2}$ on the odd reflection
sector. At $L=\log2$ the primary Arb certificate gives minimum certified
Weyl--Schur margins
\[
6.116350109501\times10^{-4}
\quad\text{and}\quad
5.261670009275\times10^{-2}
\]
in the even and odd sectors, from which~\eqref{eq:thmB-coercive} follows by
Lemma~\ref{lem:block-schur}. A clean-room reconstruction gives
$6.116514636999\times10^{-4}$ and $5.261670255041\times10^{-2}$.
\end{thmB}

The analytic step from Theorem~B to Theorem~A is exact.

\begin{thmC}[shift-energy identity]
For $0<\omega\le\tfrac12$ and $0<L\le\log2$,
\begin{equation}\label{eq:energy-identity}
I-V_{\omega,L}^{*}V_{\omega,L}
=2\int_0^\omega V_{s,L}^{*}Q_{s,L}V_{s,L}\dd s,
\end{equation}
where the integral is an improper strong-operator integral. In particular
$I-V_{\omega,L}^{*}V_{\omega,L}$ is positive definite. Moreover $V_{\omega,L}$
is compact, so positive definiteness excludes norm one and gives strict
contractivity.
\end{thmC}

The energy identity supplies a canonical positive factor. Define
\begin{equation}
(\Gcal_{\omega,L}f)(s)=\sqrt2\,Q_{s,L}^{1/2}V_{s,L}f,\qquad 0<s<\omega .
\end{equation}
Then
\begin{equation}
I-V_{\omega,L}^{*}V_{\omega,L}=\Gcal_{\omega,L}^{*}\Gcal_{\omega,L}.
\end{equation}
The shift $s$ is therefore a radial coordinate and $Q_{s,L}$ is the local
positive energy density accumulated by the boundary defect.

\subsection{Organization of the argument}

The generator certificate and the shift-energy identity are the computational
and analytic halves of one contractivity theorem. Their combination identifies
compact-window Weil positivity as the local dissipation density of the shifted
causal system. Sections~2--3 set up Suzuki's system and the completed-gamma
transfer; Section~4 derives the off-center generator and its bounded-kernel
decomposition; Sections~5--7 contain the certificate; Section~8 proves the
energy identity; Sections~9--12 discuss consequences, context, superseded
routes, and open problems; Section~13 records the verification status. The
fractional Sonine realization, compact Mosco limit, and older direct-transfer
certificates are retained as secondary structure and cross-checks.

\subsection{What is not claimed}

The scope is deliberately limited.
\begin{enumerate}[nosep]
\item The theorem does not prove RH.
\item It gives no new zero-free region. It verifies a necessary finite-section
  condition for the open shifted zero-free assertions, but a prime-free
  finite-depth compression cannot determine global innerness.
\item No arithmetic layer $n\ge2$ is controlled. The first new layer appears
  exactly at depth $\log2$.
\item The positive radial factor is a rigorous one-dimensional bulk in the
  sense of a transfer-system energy identity. It is not claimed to be an
  ordinary AdS/CFT dual.
\item The proof is computer-assisted. The primary special-function enclosure
  is formal Arb arithmetic; the independent implementation is a strong audit
  rather than a second formal special-function interval library.
\end{enumerate}

% =====================================================================
\section{Suzuki's arithmetic Hankel system}
% =====================================================================

\subsection{Prime--gamma layer decomposition}

Define the arithmetic coefficients~\cite[eq.~(2.1)]{Suzuki2012}
\begin{equation}
c_\omega(n)=n^\omega\sum_{d\mid n}\frac{\mu(d)}{d^{2\omega}}
=n^\omega\prod_{p\mid n}\bigl(1-p^{-2\omega}\bigr),
\qquad
\sum_{n\ge1}c_\omega(n)n^{-w}=\frac{\zeta(w-\omega)}{\zeta(w+\omega)} .
\end{equation}
Suzuki constructs an explicit archimedean function $g_\omega$ supported in
$(0,1)$ and the arithmetic kernel~\cite[eq.~(2.3)]{Suzuki2012}
\begin{equation}\label{eq:h-def}
h_\omega(x)=\frac1x\sum_{n\le x}c_\omega(n)\,g_\omega(n/x),\qquad x>1,
\end{equation}
with $h_\omega=0$ on $(0,1)$, and proves~\cite[Prop.~2.1]{Suzuki2012} that
$\int_0^\infty h_\omega(x)x^{1/2+iz}\dd x/x=\Theta_\omega(z)$ for
$\Im z>\tfrac12+\omega$. The multiplicative Hankel operator
is~\cite[eq.~(2.5)]{Suzuki2012}
\begin{equation}
(\Hsf_\omega f)(x)=\int_0^\infty h_\omega(xy)f(y)\dd y .
\end{equation}
Introduce logarithmic coordinates
\begin{equation}
k_\omega(t)=e^{t/2}h_\omega(e^{t}),\qquad
\kappa_\omega(t)=e^{-t/2}g_\omega(e^{-t})\mathbf 1_{t>0}.
\end{equation}
Then~\eqref{eq:h-def} becomes
\begin{equation}\label{eq:layer}
k_\omega(t)=\sum_{n\ge1}\frac{c_\omega(n)}{\sqrt n}\,\kappa_\omega(t-\log n).
\end{equation}
Every integer contributes a delayed copy of the same archimedean impulse. A
finite sum-coordinate interval of length $L<\log2$ sees only $n=1$. At
$L=\log2$ the $n=2$ contribution lies on a measure-zero boundary and does not
enter the quadratic form.

\subsection{Finite-volume Hankel operator}

On the centered interval $I_L=(-L/2,L/2)$ the first-slab Hankel operator is
\begin{equation}
(H^{\mathrm c}_{\omega,L}f)(x)=\int_{I_L}k_\omega(x+y)f(y)\dd y .
\end{equation}
Because $L\le\log2$, the layer formula~\eqref{eq:layer} gives
$k_\omega=\kappa_\omega$ almost everywhere on the relevant sum-coordinate
range. After shifting $I_L$ to $(0,L)$ the same operator is
\begin{equation}
(H_{\omega,L}f)(x)=\int_0^L\kappa_\omega(x+y-L)f(y)\dd y .
\end{equation}
In Suzuki's multiplicative normalization this is the compression
\begin{equation}
\Hsf_{\omega,a}=P_a\Hsf_\omega P_a\colon L^2((0,a),dx)\to L^2((0,a),dx)
\end{equation}
of~\cite[eq.~(2.6)]{Suzuki2012}. Since $h_\omega(xy)=0$ for $xy<1$, the
compression acts on $(1/a,a)$, which in logarithmic coordinates is the
centered interval of total length $L=2\log a$.

Define
\begin{equation}
(V_{\omega,L}f)(x)=\int_0^x\kappa_\omega(x-y)f(y)\dd y,
\qquad
(R_Lf)(x)=f(L-x).
\end{equation}
A change of variables gives the exact identities
\begin{equation}\label{eq:volterra-identities}
H_{\omega,L}=V_{\omega,L}R_L=R_LV_{\omega,L}^{*},
\qquad
R_LV_{\omega,L}R_L=V_{\omega,L}^{*},
\qquad
H_{\omega,L}^{2}=V_{\omega,L}V_{\omega,L}^{*}.
\end{equation}
Hence $H_{\omega,L}$ is self-adjoint and $\norm{H_{\omega,L}}=\norm{V_{\omega,L}}$.
It follows that
\begin{equation}
I\pm H_{\omega,L}\succ0\quad\Longleftrightarrow\quad\norm{V_{\omega,L}}<1 .
\end{equation}
This is the first-slab reflection-positivity criterion.

\subsection{Local positivity versus global innerness}

If $\Theta_\omega$ is meromorphic inner, its Hankel transform is an isometry
\cite[Lemma~4.1]{Suzuki2012} and every finite section is contractive.
Conversely, compatible contractivity at all depths recovers the global
Hardy-space mapping property in Suzuki's criterion~\cite[Thm.~2.2]{Suzuki2012}.
Positivity at one finite depth is only a local statement. In particular,
Theorem~A has no direct zero-free consequence.

% =====================================================================
\section{Completed-gamma transfer and fractional structure}
% =====================================================================

\subsection{Transfer function}

Put
\begin{equation}
\Ginf(s)=\tfrac12 s(s-1)\pi^{-s/2}\Gamma(s/2),
\end{equation}
so that $\xi=\Ginf\zeta$; this is the factor written $\gamma(s)$
in~\cite{Suzuki2012}. Suzuki shows~\cite[Sect.~3.1]{Suzuki2012} that
\begin{equation}\label{eq:suzuki-mellin}
\int_0^\infty g_\omega(x)\,x^{s}\,\frac{\dd x}{x}
=\frac{\Ginf(s-\omega)}{\Ginf(s+\omega)},
\qquad \Re s>\max(\omega-2,\,1-\omega).
\end{equation}
With $\kappa_\omega(t)=e^{-t/2}g_\omega(e^{-t})$ and $s=p+\tfrac12$ this is
the Laplace transform of the first-slab impulse,
\begin{equation}\label{eq:K}
K_\omega(p)=\int_0^\infty\kappa_\omega(t)e^{-pt}\dd t
=\frac{\Ginf(\tfrac12+p-\omega)}{\Ginf(\tfrac12+p+\omega)} .
\end{equation}
For $0<\omega<\tfrac12$ write
\begin{equation}
\alpha=\tfrac12-\omega,\qquad\beta=\tfrac12+\omega .
\end{equation}
Then
\begin{equation}\label{eq:K-factored}
K_\omega(p)=\pi^{\omega}\,\frac{p+\alpha}{p-\alpha}\,\frac{p-\beta}{p+\beta}\,
\frac{\Gamma\bigl((p+\alpha)/2\bigr)}{\Gamma\bigl((p+\beta)/2\bigr)} .
\end{equation}
The pole at $p=\alpha$ obstructs a naive stationary bounded-real argument. It
causes no difficulty on a finite horizon and disappears from the bounded
generator difference derived in Section~4.

\subsection{Explicit local impulse}

Let $q(t)=1-e^{-2t}$. For $0<\omega\le\tfrac12$ the first-slab impulse can be
written
\begin{equation}\label{eq:kappa-explicit}
\kappa_\omega(t)=\pi^\omega\Bigl[
\frac{2e^{-\alpha t}q(t)^{\omega-1}}{\Gamma(\omega)}
-\frac{4\alpha\omega\,e^{\alpha t}}{\Gamma(\omega)}B_{q(t)}(\omega,\alpha)
-\frac{4\beta\omega\,e^{-\alpha t}q(t)^{\omega}}{\Gamma(\omega+1)}
\Bigr],
\end{equation}
where $B_q(a,b)=\int_0^q u^{a-1}(1-u)^{b-1}\dd u$ and the middle term is
interpreted as zero at $\alpha=0$. Its Laplace transform is~\eqref{eq:K};
by uniqueness of Laplace transforms, \eqref{eq:kappa-explicit} agrees with
$e^{-t/2}g_\omega(e^{-t})$ almost everywhere. Near the diagonal,
$\kappa_\omega(t)=O(t^{\omega-1})$ as $t\downarrow0$, so
$\kappa_\omega\in L^1(0,L)$ for every positive shift.

\subsection{Diffusive and Sonine factors}

The fractional gamma ratio
\begin{equation}
F_\omega(p)=\frac{\Gamma((p+\alpha)/2)}{\Gamma((p+\beta)/2)}
\end{equation}
has impulse
\begin{equation}
a_\omega(t)=\frac{2}{\Gamma(\omega)}e^{-\alpha t}(1-e^{-2t})^{\omega-1}
=\sum_{n=0}^\infty d_{\omega,n}e^{-\lambda_{\omega,n}t},
\qquad
\lambda_{\omega,n}=\alpha+2n,\quad
d_{\omega,n}=\frac{2(1-\omega)_n}{\Gamma(\omega)\,n!}>0 .
\end{equation}
Define the complementary kernel
\begin{equation}
b_\omega(t)=\frac{2}{\Gamma(1-\omega)}e^{-\beta t}(1-e^{-2t})^{-\omega}.
\end{equation}
The gamma recurrence gives the Sonine identity
\begin{equation}
a_\omega*b_\omega=2e^{-\alpha t},
\qquad\text{equivalently}\qquad
\tfrac12(\partial_t+\alpha)\mathcal B_\omega\mathcal A_\omega=I
\end{equation}
on the zero-initial-value core. This places the first-slab system in the
general Sonine-kernel framework~\cite{Luchko2021}. It also gives the rational
collapse
\begin{equation}
F_{1-\omega}(p+1)K_\omega(p)=\frac{2\pi^\omega(p-\beta)}{(p+\beta)(p-\alpha)} .
\end{equation}
These identities motivate a distributed-state realization, but they are not
needed for the proof of Theorem~A.

% =====================================================================
\section{The off-center Weil generator}
% =====================================================================

\subsection{Shift derivative}

Set $G(p)=\Ginf(\tfrac12+p)$ and $\ell(p)=\log G(p)$. Then
$K_\omega(p)=G(p-\omega)/G(p+\omega)$ and
\begin{equation}\label{eq:dK}
\partial_\omega K_\omega(p)=-a_\omega(p)K_\omega(p),
\qquad
a_\omega(p)=\ell'(p-\omega)+\ell'(p+\omega).
\end{equation}
In terms of $\alpha=\tfrac12-\omega$ and $\beta=\tfrac12+\omega$, standard
gamma and digamma identities~\cite{DLMF} give
\begin{equation}\label{eq:a-mult}
a_\omega(p)=
\frac1{p+\alpha}+\frac1{p-\alpha}+\frac1{p+\beta}+\frac1{p-\beta}
-\log\pi
+\tfrac12\psi\Bigl(\frac{p+\alpha}2\Bigr)+\tfrac12\psi\Bigl(\frac{p+\beta}2\Bigr).
\end{equation}
Let $A_{\omega,L}$ be the maximal finite-horizon causal operator with transfer
$a_\omega$. On a smooth zero-initial-value core,
\begin{equation}
\partial_\omega V_{\omega,L}=-A_{\omega,L}V_{\omega,L}.
\end{equation}
Define the self-adjoint generator form
\begin{equation}\label{eq:Qdef}
Q_{\omega,L}=\tfrac12\bigl(A_{\omega,L}+A_{\omega,L}^{*}\bigr).
\end{equation}
At $\omega=0$ this is the central prime-free compact-window Weil form.

\subsection{Fourier-plus-residue representation}

Let $f$ be supported on $I_L=(-L/2,L/2)$, with Fourier transform
$\widehat f(\tau)=\int_{I_L}f(x)e^{-i\tau x}\dd x$. Define
\begin{equation}\label{eq:w}
w_\omega(\tau)=\tfrac12\Re\psi\Bigl(\frac{\alpha+i\tau}2\Bigr)
+\tfrac12\Re\psi\Bigl(\frac{\beta+i\tau}2\Bigr)-\log\pi,
\end{equation}
and the evaluation functionals
\begin{equation}
\Cc_c(f)=\int_{I_L}f(x)\cosh(cx)\dd x,\qquad
\Ss_c(f)=\int_{I_L}f(x)\sinh(cx)\dd x .
\end{equation}
For $0\le\omega<\tfrac12$, contour displacement gives
\begin{equation}\label{eq:Q-fourier}
Q_{\omega,L}[f]=\frac1{2\pi}\int_\R w_\omega(\tau)\abs{\widehat f(\tau)}^2\dd\tau
+\abs{\Cc_\alpha(f)}^2+\abs{\Cc_\beta(f)}^2
-\abs{\Ss_\alpha(f)}^2-\abs{\Ss_\beta(f)}^2 .
\end{equation}
Even and odd reflection sectors therefore decouple. The endpoint needs a
distributional correction and is recorded explicitly.

\begin{lemma}[right-endpoint boundary term]\label{lem:endpoint}
Let $w^{\pt}_{1/2}$ be the pointwise limit of $w_\omega$ as
$\omega\uparrow\tfrac12$ away from $\tau=0$, continuously extended at zero.
Then
\begin{equation}\label{eq:Q-endpoint}
Q_{1/2,L}[f]=\frac1{2\pi}\int_\R w^{\pt}_{1/2}(\tau)\abs{\widehat f(\tau)}^2\dd\tau
+\tfrac12\abs{\Cc_0(f)}^2+\abs{\Cc_1(f)}^2-\abs{\Ss_1(f)}^2 .
\end{equation}
\end{lemma}

\begin{proof}
Since $\psi(z)=-1/z+O(1)$ near $z=0$,
\[
\tfrac12\Re\psi\Bigl(\frac{\alpha+i\tau}2\Bigr)
=-\frac{\alpha}{\alpha^2+\tau^2}+\tfrac12\Re\psi\Bigl(\frac{i\tau}2\Bigr)+o(1),
\]
and $\int_\R\alpha(\alpha^2+\tau^2)^{-1}\dd\tau=\pi$, so
$\tfrac12\Re\psi((\alpha+i\tau)/2)\to\tfrac12\Re\psi(i\tau/2)-\pi\delta_0$
in distributions as $\alpha\downarrow0$. Since
$\widehat f(0)=\Cc_0(f)=\int f$, the delta term subtracts
$\tfrac12\abs{\Cc_0(f)}^2$ from the pointwise-limit formula, while the residue
terms $\abs{\Cc_\alpha}^2-\abs{\Ss_\alpha}^2$ tend to $\abs{\Cc_0}^2$. This is
exactly the boundary contribution carried automatically by the bounded-kernel
representation of Proposition~\ref{prop:bounded-difference} below, whose
uniform convergence as $\omega\uparrow\tfrac12$ justifies the limit.
\end{proof}

\subsection{Exact bounded perturbation}

Define
\begin{equation}\label{eq:R}
R(t)=e^{t/2}-\frac{e^{-5t/2}}{1-e^{-2t}} .
\end{equation}

\begin{proposition}[bounded shift difference]\label{prop:bounded-difference}
For $0\le\omega\le\tfrac12$,
\begin{equation}\label{eq:Q-decomp}
Q_{\omega,L}=Q_{0,L}+C_{\omega,L},
\end{equation}
where $C_{\omega,L}$ is the self-adjoint integral operator on $I_L$ with kernel
\begin{equation}\label{eq:c-kernel}
c_\omega(x,y)=\bigl(\cosh(\omega\abs{x-y})-1\bigr)R(\abs{x-y}) .
\end{equation}
\end{proposition}

\begin{proof}
First take $0\le\omega<\tfrac12$ and $\Re p$ sufficiently large. The
convergent difference must be formed before its summands are separated. From
$\psi(z)=-\gamma+\int_0^\infty(e^{-u}-e^{-zu})(1-e^{-u})^{-1}\dd u$ with
$u=2t$,
\begin{equation}
\tfrac12\psi\Bigl(\frac{p+c}2\Bigr)=-\frac\gamma2+\int_0^\infty
\frac{e^{-2t}-e^{-(p+c)t}}{1-e^{-2t}}\dd t .
\end{equation}
In the sum with $c=\alpha,\beta$, minus twice the $c=\tfrac12$ term, the
constants and the individually divergent $e^{-2t}$ terms cancel. The remaining
integrand is integrable at zero and equals
\begin{equation}
-2e^{-pt}\,\frac{e^{-t/2}}{1-e^{-2t}}\bigl(\cosh(\omega t)-1\bigr).
\end{equation}
Independently, the four rational terms in $a_\omega-a_0$ have inverse Laplace
kernel
\begin{equation}
e^{-\alpha t}+e^{\alpha t}+e^{-\beta t}+e^{\beta t}-4\cosh(t/2)
=4\cosh(t/2)\bigl(\cosh(\omega t)-1\bigr),
\end{equation}
using $\cosh\alpha t+\cosh\beta t=2\cosh(t/2)\cosh(\omega t)$. Adding the two
contributions and using
\begin{equation}
2\cosh(t/2)-\frac{e^{-t/2}}{1-e^{-2t}}
=e^{t/2}-\frac{e^{-5t/2}}{1-e^{-2t}}=R(t)
\end{equation}
yields the multiplier identity
\begin{equation}\label{eq:multiplier-identity}
a_\omega(p)-a_0(p)=\int_0^\infty e^{-pt}\,2\bigl(\cosh(\omega t)-1\bigr)R(t)\dd t .
\end{equation}
The causal convolution with this kernel has the kernel
$2(\cosh(\omega(x-y))-1)R(x-y)\mathbf 1_{x>y}$; symmetrizing according
to~\eqref{eq:Qdef} removes the factor two and gives~\eqref{eq:c-kernel}.
Finally, $c_\omega$ converges uniformly on the compact square $I_L^2$ as
$\omega\uparrow\tfrac12$ (the factor $\cosh(\omega t)-1$ is entire in
$\omega$ and $R$ is fixed); hence the form identity extends to the endpoint
by bounded-operator convergence, consistently with Lemma~\ref{lem:endpoint}.
\end{proof}

\subsection{Common logarithmic form domain}

As $t\downarrow0$, $R(t)=-\tfrac1{2t}+O(1)$ while
$\cosh(\omega t)-1=\tfrac12\omega^2t^2+O(t^4)$. Thus $c_\omega(x,y)=O(\abs{x-y})$
at the diagonal. Hence $C_{\omega,L}$ is Hilbert--Schmidt, and therefore
bounded and compact. All $Q_{\omega,L}$ are closed forms on the common domain
\begin{equation}\label{eq:Dlog}
\Dlog(I_L)=\Bigl\{f\in L^2(I_L):\ \int_\R\log(2+\abs\tau)\abs{\widehat f(\tau)}^2\dd\tau<\infty\Bigr\}.
\end{equation}
This also removes the apparent parameter singularity at $\omega=\tfrac12$.

% =====================================================================
\section{The central prime-free Weil form}
% =====================================================================

\subsection{Central form}

At $\omega=0$, $\alpha=\beta=\tfrac12$ and the generator becomes
\begin{equation}\label{eq:Q0}
Q_{0,L}[f]=\frac1{2\pi}\int_\R\Bigl[\Re\psi\Bigl(\frac14+\frac{i\tau}2\Bigr)-\log\pi\Bigr]
\abs{\widehat f(\tau)}^2\dd\tau
+2\abs{\Cc_{1/2}(f)}^2-2\abs{\Ss_{1/2}(f)}^2 .
\end{equation}
This is the prime-free archimedean Weil form restricted to the compact window
$I_L$. It is diagonal with respect to reflection parity. The logarithmic
growth of the digamma multiplier makes $Q_{0,L}$ unbounded above on $L^2$,
but it is lower bounded and closed on $\Dlog$. The residue is positive in the
even sector and negative in the odd sector, so both sectors must be checked.

\subsection{Finite-frequency reduction}

At the maximal window $L=\log2$ set $T_0=\tfrac12\log2$. Let $\phi_n$ be the
normalized Legendre basis on $(-T_0,T_0)$. Its Fourier and pole-evaluation
coefficients are
\begin{equation}
\widehat\phi_n(\tau)=\sqrt{2T_0(2n+1)}\,(-i)^n j_n(T_0\tau),
\qquad
p_n=\sqrt{2T_0(2n+1)}\,i_n(T_0/2),
\end{equation}
where $j_n$ and $i_n$ are the spherical Bessel functions of the first kind and
their modified counterparts. Choose the frequency cutoff $\Lambda=30$ and the
high-frequency floor
\begin{equation}
\beta_*=\log\frac{\Lambda}{2\pi}-\frac1\Lambda=1.5299869819194765\ldots
\end{equation}
The central matrix is the sum of $\beta_*I$, a compact frequency integral
over $[0,30]$, and the parity-dependent rank-one residue. The complement is
controlled by spherical-Bessel majorants.

\subsection{Central interval certificate}

A sixteen-mode block in each parity sector was evaluated using 80-decimal Arb
arithmetic and 48-point Gauss--Legendre quadrature on each of thirty unit
frequency panels. All quadrature nodes, weights, digamma values, and Bessel
values were ball-enclosed. An entrywise radius $10^{-10}$ was added, larger
than the analytic quadrature remainder. Interval $LDL^{\mathsf T}$ succeeds
after subtracting $10^{-3}I$ in the even sector and $5\times10^{-2}I$ in the
odd sector. The head--tail couplings are bounded by
\begin{equation}\label{eq:central-couplings}
1.4084884\times10^{-9}\quad\text{and}\quad2.2191643\times10^{-10}
\end{equation}
in the even and odd sectors, and the tail deviations below $\beta_*$ by
$9.7777\times10^{-24}$ and $2.4269\times10^{-25}$. Consequently, by the Weyl
inequality applied to the off-diagonal blocks,
\begin{equation}
Q_{0,\log2}[f]\ \ge\
\begin{cases}
9.9999859\times10^{-4}\,\norm f_2^2, & f\ \text{even},\\[1mm]
4.99999998\times10^{-2}\,\norm f_2^2, & f\ \text{odd}.
\end{cases}
\end{equation}
The computation is independent of zeta zeros and RH. At $\omega=0$ it is a
certified reproduction of the classical first-window Weil positivity result,
not a claim of a new central theorem; see also the screw-form
formulation~\cite{Suzuki2026screw} and recent compact-window and Galerkin
investigations~\cite{ConnesConsani2020,Zhu2026,Kim2026,Groskin2026}. The new
part here is the parameter-uniform off-center deformation and its use as the
density in the shift-energy identity.

There is no conflict between the present first-slab boundary and the larger
certified windows in~\cite{Zhu2026}. That paper writes
$\supp f\subset[-\LZ,\LZ]$ and its prime sum contains $\log n<2\LZ$. Our
centered interval has total length $L$, so the conventions are related by
$L=2\LZ$. Thus the classical prime-free cutoff $2\LZ\le\log2$ is exactly our
$L\le\log2$, while the certificate in~\cite{Zhu2026} at $\LZ=0.8$ is
prime-active and concerns the full Weil form, not the archimedean
continuation studied here.

% =====================================================================
\section{Entire shift expansion}
% =====================================================================

\subsection{Analytic profile}

Put $S(t)=tR(t)$. Then
\begin{equation}
S(t)=te^{t/2}-\tfrac12e^{-3t/2}\frac{t}{\sinh t},
\end{equation}
which is analytic in a disk containing $[0,\log2]$. Since
$\cosh(\omega t)-1=\sum_{k\ge1}\omega^{2k}t^{2k}/(2k)!$, the perturbation is
even and entire in $\omega$:
\begin{equation}
C_{\omega,L}=\sum_{k\ge1}\omega^{2k}T_{k,L},
\qquad
(T_{k,L})(x,y)=\frac{\abs{x-y}^{2k}R(\abs{x-y})}{(2k)!}.
\end{equation}

\subsection{Exact Legendre correlations}

Let $m,n$ have the same parity. For $0\le u\le1$ define
\begin{equation}
C_{mn}(u)=\tfrac12\int_{-1}^{1-2u}\bigl[P_m(r+2u)P_n(r)+P_n(r+2u)P_m(r)\bigr]\dd r ,
\end{equation}
a polynomial with rational coefficients. If $S(t)=\sum_{d\ge0}s_dt^d$, then
\begin{equation}\label{eq:Tk-entries}
(T_{k,L})_{mn}=\frac{\sqrt{(2m+1)(2n+1)}}{(2k)!}
\sum_j[u^j]C_{mn}(u)\sum_d\frac{s_dL^{d+2k}}{d+2k+j} .
\end{equation}
Thus every finite perturbation matrix is reduced to exact rational Legendre
algebra and a scalar analytic series.

\subsection{Kernel and Taylor remainder bounds}\label{sec:kernel-bound}

The profile $R$ changes sign once on $(0,\log2)$, at
\begin{equation}
t_0=\log\rho=0.2811995743229618465\ldots,\qquad\rho^3=\rho+1,
\end{equation}
where $\rho=1.3247179572\ldots$ is the plastic number. We locate the actual
maximum row sum rather than doubling an endpoint row. Put
\begin{equation}
h(t)=\bigl(\cosh(t/2)-1\bigr)R(t),\qquad q=e^{t/2}.
\end{equation}
Then
\begin{equation}
h(t)=\frac{(q-1)(q^6-q^2-1)}{2q^2(q+1)(q^2+1)},
\end{equation}
and $\dd h/\dd t$ has the sign of
\begin{equation}
N(q)=q^{10}+q^9+q^8+q^7-q^6-q^5+q^4-2q^3-q^2-q-1 .
\end{equation}
Indeed,
\begin{equation}
N'(q)=(10q^9-6q^5)+(9q^8-5q^4)+(8q^7-6q^2)+(7q^6-2q)+(4q^3-1)>0
\end{equation}
for $q\ge1$, since each bracket is nonnegative at $q=1$ and increasing. Also
$N(1)=-2<0$ and $N(\sqrt\rho)=3\rho^2-1>0$. Thus $h$ has a unique minimum
$t_m$ before $t_0$ and is increasing after $t_m$. The Arb check encloses
\begin{equation}
\begin{aligned}
t_m&=0.1416454811943886\ldots,\\
h(t_m)&=-0.0044442714714077\ldots,\\
h(\log2-t_0)&=0.0126146638556245\ldots,
\end{aligned}
\end{equation}
with $h(\log2-t_0)+h(t_m)>0.00817039$.

\begin{lemma}[endpoint row maximum]\label{lem:row-max}
For $0<L\le\log2$ and $0\le\omega\le\tfrac12$,
\begin{equation}\label{eq:C-norm}
\norm{C_{\omega,L}}\le\int_0^{\log2}\abs{h(t)}\dd t
\le0.012293050401142320<0.012294 .
\end{equation}
\end{lemma}

\begin{proof}
Take first $L=\log2$ and $\omega=\tfrac12$, so that $\abs{c_\omega(x,y)}=\abs{h(\abs{x-y})}$.
For a row indexed by $x\in I_L$, let its left and right available lengths be
$a,b$, with $a+b=L$, and assume $a\ge b$. Its absolute row integral is
$F(a)+F(b)$, where $F(r)=\int_0^r\abs{h(t)}\dd t$, and
$\frac{\dd}{\dd b}[F(L-b)+F(b)]=\abs{h(b)}-\abs{h(a)}$. Note that
$a\ge L/2=0.3466\ldots>t_0=0.2812\ldots$, so $a\le t_0$ cannot occur and
$h(a)>0$ always. If $b\ge t_0$, monotonicity of $h$ on $[t_m,\infty)$ gives
$h(a)\ge h(b)\ge0$. If $b\le t_0$, then $h(b)\ge h(t_m)$ and
$h(a)\ge h(\log2-t_0)$ since $a\ge\log2-t_0>t_m$, so the displayed interval
check gives $h(a)+h(b)>0$, again implying $\abs{h(a)}\ge\abs{h(b)}$. Hence the
absolute row integral is nonincreasing in $b$ on $[0,L/2]$ and is maximal at
an endpoint row, where it equals $\int_0^{\log2}\abs h$. The Schur test for
the symmetric kernel gives the first inequality; sign-split Arb integration
gives the second. The bound is uniform for $0\le\omega\le\tfrac12$ and
$0<L\le\log2$, since $\abs{c_\omega}$ is pointwise increasing in $\omega$ and
restriction to a shorter centered interval can only shorten rows.
\end{proof}

Truncating the shift expansion after $\omega^8T_4$ gives
\begin{equation}\label{eq:taylor-remainder}
\sup_{0\le\omega\le1/2}\Bigl\|C_{\omega,L}-\sum_{k=1}^{4}\omega^{2k}T_{k,L}\Bigr\|
<1.54\times10^{-12}.
\end{equation}
The primary Arb enclosure, obtained from a conservative absolute-profile Cauchy
bound rather than a sharp sign split, is $1.532967865231399\times10^{-12}$.

% =====================================================================
\section{Computer-assisted generator positivity}
% =====================================================================

\subsection{Parameter boxes}

The interval $0\le\omega\le\tfrac12$ is divided into $256$ equal boxes. For
each box and parity sector the calculation forms
\begin{equation}
Q_0+[\omega^2]T_1+[\omega^4]T_2+[\omega^6]T_3+[\omega^8]T_4,
\end{equation}
where every occurrence of $\omega$ is represented by the same interval
parameter. This retains correlation between all matrix entries.

The sixteen-mode head is combined with the infinite Legendre tail. The central
tail lower bound is reduced by the full perturbation norm:
\begin{equation}\label{eq:dtail}
d_{\mathrm{tail}}\ \ge\ \beta_*-\delta_{\mathrm{central}}-0.012294
\ >\ 1.5176929819194764 .
\end{equation}
Here $\delta_{\mathrm{central}}$ is not set to zero: its certified values are
$9.7777\times10^{-24}$ in the even sector and $2.4269\times10^{-25}$ in the
odd sector, and therefore disappear at the displayed decimal precision. The
head--tail coupling is bounded by
\begin{equation}\label{eq:kappa}
\kappa=\kappa_{\mathrm{central}}+0.012294,
\end{equation}
with $\kappa_{\mathrm{central}}$ from~\eqref{eq:central-couplings}. The
resulting worst Schur loss $\kappa^2/d_{\mathrm{tail}}$ is
$9.95869865858903\times10^{-5}$ (even) and $9.95869673623490\times10^{-5}$
(odd).

\subsection{Primary Arb result}

For every box, interval $LDL^{\mathsf T}$ proves positivity of the shifted
head after subtracting the Schur loss and Taylor remainder. All $512$
parity-box problems pass. The minimum certified quantities are as follows.

\begin{center}
\begin{tabular}{lrr}
\toprule
Quantity & Even & Odd\\
\midrule
Weyl--Schur margin $m$ & $6.116350109501\times10^{-4}$ & $5.261670009275\times10^{-2}$\\
tail lower bound $d_{\mathrm{tail}}$ & $1.51769298191948$ & $1.51769298191948$\\
Schur loss $\kappa^2/d_{\mathrm{tail}}$ & $9.95869865858903\times10^{-5}$ & $9.95869673623490\times10^{-5}$\\
\bottomrule
\end{tabular}
\end{center}

Across both sectors the smallest interval $LDL^{\mathsf T}$ pivot is
$2.326811914886\times10^{-2}$. The margin $m$ is, by construction, a certified
lower bound for $\lambda_{\min}(\text{head})-\kappa^2/d_{\mathrm{tail}}$ on
each box, with the interval widening of the box, the entrywise radii, and the
Taylor remainder~\eqref{eq:taylor-remainder} already subtracted.

\subsection{From the certified margin to coercivity on the form domain}

The certificate is a statement about a $16\times16$ head, an infinite tail,
and their coupling. The following elementary lemma converts it into a
coercivity bound for the full form on $\Dlog$, which is the form in which
Section~8.6 consumes it.

\begin{lemma}[block-Schur coercivity]\label{lem:block-schur}
Let $Q$ be a closed lower-bounded form with domain $\mathcal D\subset\Hcal$, let
$\Hcal=\Hcal_h\oplus\Hcal_t$ with $\Hcal_h$ finite-dimensional and
$\Hcal_h\subset\mathcal D$, and suppose that for $f_h\in\Hcal_h$ and
$f_t\in\mathcal D\cap\Hcal_t$,
\[
Q[f_h]\ge\lambda\norm{f_h}^2,\qquad
Q[f_t]\ge d\,\norm{f_t}^2,\qquad
\abs{Q[f_h,f_t]}\le\kappa\norm{f_h}\norm{f_t},
\]
with $d>0$. Then $Q[f]\ge c\norm f^2$ for all $f\in\mathcal D$ and every
$0\le c<\min(\lambda,d)$ with $(\lambda-c)(d-c)\ge\kappa^2$. In particular, if
$\lambda\ge m+\kappa^2/d$ with $0<m<d$, then $Q\succeq c\,I$ for
\begin{equation}\label{eq:c-from-m}
c=m-\frac{\kappa^2m}{d(d-m)} .
\end{equation}
\end{lemma}

\begin{proof}
For $f=f_h+f_t$ with $f_h\in\Hcal_h$ and $f_t=f-f_h\in\mathcal D\cap\Hcal_t$,
\[
Q[f]=Q[f_h]+2\Re Q[f_h,f_t]+Q[f_t]
\ge\lambda\norm{f_h}^2-2\kappa\norm{f_h}\norm{f_t}+d\norm{f_t}^2,
\]
and the quadratic form on the right dominates $c(\norm{f_h}^2+\norm{f_t}^2)$
exactly when the matrix $\begin{psmallmatrix}\lambda-c&-\kappa\\-\kappa&d-c\end{psmallmatrix}$
is positive semidefinite, i.e.\ when $\lambda\ge c$, $d\ge c$, and
$(\lambda-c)(d-c)\ge\kappa^2$. For the second statement put $c=m-\delta$ with
$\delta=\kappa^2m/(d(d-m))$; then
$(\lambda-c)(d-c)\ge(\kappa^2/d+\delta)(d-m+\delta)\ge\kappa^2+\delta(d-m)-\kappa^2m/d=\kappa^2$.
\end{proof}

\begin{remark}\label{rem:not-m}
The naive reading ``$Q\succeq m\,I$'' is \emph{not} implied by the
certificate: the Schur complement of the tail at level $c$ is
$\kappa^2/(d-c)$, not $\kappa^2/d$, and the difference costs
$\kappa^2m/(d(d-m))$. Because $d_{\mathrm{tail}}\approx1.5177\gg m$ the
correction is small---$4.0\times10^{-8}$ in the even sector and
$3.6\times10^{-6}$ in the odd sector---but it is the reason Theorem~B states
$c_*=6.1159\times10^{-4}$ rather than the certified margin itself. With the
data of Section~7.2, \eqref{eq:c-from-m} gives
$c_{\mathrm{even}}=6.11594\ldots\times10^{-4}$ and
$c_{\mathrm{odd}}=5.26131\ldots\times10^{-2}$.
\end{remark}

\begin{proof}[Proof of Theorem~B]
Fix a parity sector and $L=\log2$. Take $\Hcal_h$ to be the span of the
sixteen head Legendre modes of that parity, $\Hcal_t$ its orthogonal
complement in the sector, and $\mathcal D=\Dlog(I_{\log2})$ intersected with
the sector. For each box, the certificate gives
$\lambda\ge m+\kappa^2/d_{\mathrm{tail}}$ on the head,
$d_{\mathrm{tail}}$ on the tail by~\eqref{eq:dtail}, and $\kappa$
by~\eqref{eq:kappa}. Lemma~\ref{lem:block-schur} yields
$Q_{\omega,\log2}\succeq c_{\mathrm{even}}I$ and
$Q_{\omega,\log2}\succeq c_{\mathrm{odd}}I$ on the two sectors, for every
$\omega$ in the box; the union of the boxes is $[0,\tfrac12]$ and the two
sectors are orthogonal and $Q$-invariant, so~\eqref{eq:thmB-coercive} holds at
$L=\log2$ with $c_*=\min(c_{\mathrm{even}},c_{\mathrm{odd}})\ge6.1159\times10^{-4}$.

If $0<L'<\log2$, zero extension identifies $L^2(I_{L'})$ with a closed subspace
of $L^2(I_{\log2})$ and preserves parity. The
representations~\eqref{eq:Q-fourier}--\eqref{eq:Q-endpoint} depend on $f$ only
through $\widehat f$ and the functionals $\Cc_c(f),\Ss_c(f)$, none of which
depends on the horizon, so $Q_{\omega,L'}$ is exactly the restriction of
$Q_{\omega,\log2}$ to that subspace, and $\Dlog(I_{L'})$ is the corresponding
restriction of the form domain. Positivity with the same constant therefore
passes to every shorter centered interval. In the causal coordinates of
Section~2, the same statement says that causality makes $V_{\omega,L'}$ the
finite-horizon compression of $V_{\omega,\log2}$.
\end{proof}

\subsection{Endpoint safety}

The certificate includes $\omega=\tfrac12$. No limiting digamma evaluation is
used there. The bounded-kernel representation $Q_\omega=Q_0+C_\omega$ continues
analytically through the endpoint and carries the boundary delta contribution
that would be lost by a pointwise multiplier limit. Explicitly, that
contribution is $-\pi\delta_0$ in the multiplier, or $-\tfrac12\abs{\int f}^2$
in the quadratic form, as stated in Lemma~\ref{lem:endpoint}.

\subsection{Exact role of the computer}

The following inputs are analytic:
\begin{enumerate}[nosep]
\item the transfer identity~\eqref{eq:multiplier-identity} for $a_\omega-a_0$;
\item the bounded-kernel formula~\eqref{eq:c-kernel};
\item the common logarithmic form domain~\eqref{eq:Dlog};
\item the rational Legendre correlation formula~\eqref{eq:Tk-entries};
\item the scalar Cauchy remainder for $S$;
\item the kernel Schur bound of Lemma~\ref{lem:row-max};
\item the spherical-Bessel infinite-tail estimates;
\item the block-Schur reduction of Lemma~\ref{lem:block-schur}.
\end{enumerate}
The computer performs interval evaluation of the finite matrices and interval
$LDL^{\mathsf T}$ on the $512$ boxes. The theorem does not infer positivity
from a floating grid.

\subsection{Independent implementation audit}\label{sec:audit}

A clean-room audit implementation imports neither the primary generator code
nor the central-certificate code. Its independent paths are:
\begin{itemize}[nosep]
\item arbitrary-precision digamma and Bessel evaluation rather than Arb;
\item separately generated 16- and 48-point quadrature rules;
\item the closed binomial formula for Legendre polynomials rather than recurrence;
\item the Bernoulli expansion of $t/\sinh t$ rather than reciprocal-series inversion;
\item direct nonlinear-kernel integrals at selected entries;
\item a custom 80-digit directed-rounding \texttt{Decimal} interval class; and
\item an independent interval $LDL^{\mathsf T}$ implementation.
\end{itemize}
The maximum central quadrature discrepancy is $1.6931\times10^{-15}$. The
maximum difference between the independently reconstructed and primary
floating minimum eigenvalues is $4.3755\times10^{-13}$. The maximum
discrepancy between direct nonlinear-kernel integration and the degree-four
Taylor reconstruction is $1.5929\times10^{-13}$. All $512$ boxes pass. The
independently recovered minimum margins are
$6.116514636999\times10^{-4}$ and $5.261670255041\times10^{-2}$, and the
minimum directed-interval pivot is $2.326829938890\times10^{-2}$.

The audit assigns every central entry a radius $5\times10^{-10}$, over five
orders of magnitude larger than the observed independent quadrature
discrepancy. Its final margins are nevertheless slightly larger than the
primary margins because the two programs propagate the shift-polynomial
coefficient uncertainty differently. On the weakest even box, the audit
midpoint is $1.57\times10^{-9}$ lower but its row-sum radius is
$1.80\times10^{-8}$ smaller, producing a net margin increase of
$1.65\times10^{-8}$. On the weakest odd box, the corresponding changes are
$-2.12\times10^{-8}$ in the midpoint and $-2.36\times10^{-8}$ in the radius,
producing a net increase of $2.46\times10^{-9}$. Thus the audit's larger fixed
central radius is more than offset by its tighter directed interval
propagation of the parameter-dependent coefficients; the two final intervals
are independent enclosures and are not asserted to be nested. This is a strong
independent reproduction. Its one qualification is that the second
special-function path is arbitrary-precision rather than a second formally
validated interval special-function library. The original Arb calculation
remains the formal special-function enclosure. The endpoint-row-maximum
\emph{argument} of Lemma~\ref{lem:row-max} is inherited by the audit from the
primary implementation; only the scalar integral in~\eqref{eq:C-norm} was
recomputed independently.

% =====================================================================
\section{The shift-energy identity}
% =====================================================================

\subsection{Weighted Hardy realization}

Fix $\eta>1$ and consider $\Hcal_\eta=L^2(\R_+,e^{-2\eta t}\dd t)$. The Laplace
transform on $\Re p=\eta$ identifies this space with a Hardy boundary space.
Multiplication by $K_s$ defines a bounded causal operator $T_s$, and
\begin{equation}
V_{s,L}=P_LT_sE_L,
\end{equation}
where $E_L$ is zero extension and $P_L$ is restriction. The choice $\eta>1$
keeps the line $\Re p=\eta$ to the right of the pole $p=\alpha<\tfrac12$
of~\eqref{eq:K-factored}.

Uniform Stirling estimates on the fixed right half-plane give
\begin{equation}
K_s(\eta+i\tau)=O\bigl((2+\abs\tau)^{-s}\bigr)
\end{equation}
uniformly for $0\le s\le\tfrac12$. For every $\eps>0$ and $m=0,1,2$,
\begin{equation}
\partial_s^mK_s(\eta+i\tau)=O\bigl((1+\log(2+\abs\tau))^m(2+\abs\tau)^{-\eps}\bigr)
\end{equation}
uniformly for $\eps\le s\le\tfrac12$. Also
$a_s(\eta+i\tau)=O(1+\log(2+\abs\tau))$, so $a_sK_s$ is a bounded multiplier
for every positive $s$. Consequently $s\mapsto V_{s,L}$ is continuously
differentiable in operator norm on every compact positive-shift interval, and
\begin{equation}\label{eq:V-prime}
V_{s,L}'=-A_{s,L}V_{s,L}.
\end{equation}
Although $A_{s,L}$ is unbounded, the product $A_{s,L}V_{s,L}$ is bounded. The
weighted Hardy realization is used only to prove this adjoint-free derivative
identity. On the finite interval the weighted and unweighted $L^2$ norms are
equivalent, so the identity and its norm-differentiability transfer to
operators on unweighted $L^2(0,L)$. Every adjoint from Section~8.4 onward is
taken with respect to that unweighted inner product.

\subsection{Strong limit at zero}

For each fixed $p$ on $\Re p=\eta$, $K_s(p)\to1$ as $s\downarrow0$. The
uniform multiplier bound and dominated convergence give
\begin{equation}\label{eq:strong-limit}
V_{s,L}f\longrightarrow f\quad\text{in }L^2(0,L)
\end{equation}
for every $f\in L^2(0,L)$. This convergence is strong, not in operator norm:
the high-frequency decay $K_s(\eta+i\tau)=O(\abs\tau^{-s})$ prevents uniform
multiplier convergence to one.

\subsection{Smoothing into the common form domain}

For $0<r<s$, $\sup_\tau(1+\abs\tau)^r\abs{K_s(\eta+i\tau)}<\infty$. Choose
$r<\min(s,\tfrac12)$. The full causal multiplier maps $L^2$ into $H^r$, and
restriction followed by zero extension is bounded on $H^r$ for $r<\tfrac12$.
Since $\log(2+\abs\tau)\le C_r(1+\abs\tau)^{2r}$, we obtain
\begin{equation}
V_{s,L}L^2(0,L)\subset\Dlog(I_L).
\end{equation}
By closure from the smooth core,
\begin{equation}\label{eq:Q-pairing}
Q_{s,L}[V_{s,L}f]=\Re\langle A_{s,L}V_{s,L}f,V_{s,L}f\rangle .
\end{equation}

\subsection{Energy identity away from zero}

Let $D_{s,L}=I-V_{s,L}^{*}V_{s,L}$. For $0<\eps<\omega\le\tfrac12$,
operator-norm differentiation gives
\begin{equation}
\frac{\dd}{\dd s}\langle D_{s,L}f,g\rangle
=-\langle V_{s,L}'f,V_{s,L}g\rangle-\langle V_{s,L}f,V_{s,L}'g\rangle
=2Q_{s,L}[V_{s,L}f,V_{s,L}g].
\end{equation}
The right-hand side represents a norm-continuous bounded operator on
$[\eps,\omega]$. Therefore
\begin{equation}
D_{\omega,L}-D_{\eps,L}=2\int_\eps^\omega V_{s,L}^{*}Q_{s,L}V_{s,L}\dd s,
\end{equation}
with an operator-norm integral.

\subsection{Passage to zero}

By the strong approximate-identity result~\eqref{eq:strong-limit},
$D_{\eps,L}\to0$ strongly. Theorem~B makes the integrand nonnegative. Letting
$\eps\downarrow0$ gives
\begin{equation}
D_{\omega,L}=2\int_0^\omega V_{s,L}^{*}Q_{s,L}V_{s,L}\dd s\ \succeq\ 0,
\end{equation}
where the improper integral is the monotone strong limit of the norm integrals
over $[\eps,\omega]$. Equivalently,
\begin{equation}\label{eq:scalar-energy}
\langle D_{\omega,L}f,f\rangle=2\int_0^\omega Q_{s,L}[V_{s,L}f]\dd s .
\end{equation}

\subsection{Strictness}\label{sec:strictness}

Suppose $\langle D_{\omega,L}f,f\rangle=0$. Theorem~B is uniform in the shift
and gives, for the maximal horizon and hence for every compression,
\begin{equation}
Q_{s,L}[g]\ge c_*\norm g_2^2,\qquad c_*=6.1159\times10^{-4},
\qquad0\le s\le\tfrac12,\quad g\in\Dlog(I_L).
\end{equation}
This is a bound on the whole form domain, not only on the certified head:
Lemma~\ref{lem:block-schur} converts the certified head margin into it. The
scalar integrand in~\eqref{eq:scalar-energy} is measurable and nonnegative.
Its integral can vanish only if $Q_{s,L}[V_{s,L}f]=0$ for almost every
$s\in(0,\omega)$. The uniform bound implies $V_{s,L}f=0$ for almost every such
$s$. Choose a sequence from this full-measure set with $s_n\downarrow0$. The
strong limit $V_{s_n,L}f\to f$ then yields $f=0$. Hence $D_{\omega,L}>0$.

For fixed $s>0$, the weak singularity $\kappa_s(t)=O(t^{s-1})$ is integrable.
Truncating $\kappa_s$ away from $t=0$ gives Hilbert--Schmidt Volterra
operators, while Young's inequality shows convergence in operator norm. Thus
$V_{s,L}$ is compact. If $\norm{V_{\omega,L}}=1$, the compact positive operator
$V_{\omega,L}^{*}V_{\omega,L}$ has eigenvalue one, contradicting the positive
definiteness of $D_{\omega,L}$. Therefore $\norm{V_{\omega,L}}<1$. This proves
Theorem~C and, with~\eqref{eq:volterra-identities}, Theorem~A on
$0<\omega\le\tfrac12$.

% =====================================================================
\section{Consequences and interpretations}
% =====================================================================

\subsection{Reflection positivity}

Since $H_{\omega,L}=V_{\omega,L}R_L$ is self-adjoint and
$\norm{H_{\omega,L}}=\norm{V_{\omega,L}}<1$, we have $I\pm H_{\omega,L}\succ0$.
This is reflection positivity for the two finite-volume Hankel channels, in
the operator-positivity sense inherited from the Osterwalder--Schrader
framework~\cite{OS1973}. Shorter horizons follow either from the same proof or
by compression from $L=\log2$.

\subsection{Shift-radial Gram factor}

Let $\Kfrak_{\omega,L}=L^2\bigl((0,\omega);L^2(0,L)\bigr)$ and define
$(\Gcal_{\omega,L}f)(s)=\sqrt2\,Q_{s,L}^{1/2}V_{s,L}f$. The energy identity
gives
\begin{equation}
2\int_0^\omega\norm{Q_{s,L}^{1/2}V_{s,L}f}^2\dd s=\langle D_{\omega,L}f,f\rangle\le\norm f^2 .
\end{equation}
Hence $\Gcal_{\omega,L}$ is bounded and
$D_{\omega,L}=\Gcal_{\omega,L}^{*}\Gcal_{\omega,L}$. The boundary defect is
therefore the accumulated positive energy along the shift coordinate.

\subsection{Context: the classical shifted range \texorpdfstring{$\omega>\tfrac12$}{omega > 1/2}}\label{sec:known-range}

For $\omega>\tfrac12$ the function $\Theta_\omega$ is inner unconditionally:
$\zeta(s)\neq0$ for $\Re s\ge1$, so~\cite[Prop.~1.2]{Suzuki2012} applies with
$\omega_0=\tfrac12$. Global innerness makes $\Hsf_\omega$ an
isometry~\cite[Lemma~4.1]{Suzuki2012} and gives the non-strict estimate
$\norm{V_{\omega,L}}\le1$ for every finite section. Suzuki goes further:
for $\omega>\tfrac12$ and every $a>1$ he proves~\cite[Lemma~4.4]{Suzuki2012}
that $\norm{\Hsf_{\omega,a}}<1$, so that $1\pm\Hsf_{\omega,a}$ are invertible.
The strictness there comes from a support lemma~\cite[Lemma~4.3]{Suzuki2012}
of Paley--Wiener type---for $\omega>\tfrac12$, $\Hsf_\omega P_af$ cannot have
compact support unless it vanishes---combined with the isometry and the
compactness of $\Hsf_{\omega,a}$. In the normalization of Section~2.2 this
covers all $\omega>\tfrac12$ and all $L=2\log a>0$, at every depth, prime
layers included.

Theorem~A is the complementary statement on the critical interval. It covers
$0<\omega\le\tfrac12$, where no zero-free input is available and Suzuki's
support lemma is not stated, at the prime-free depths $L\le\log2$. At the
endpoint $\omega=\tfrac12$, where~\cite[Prop.~1.2]{Suzuki2012} gives nothing
and~\cite[Lemma~4.4]{Suzuki2012} is not stated, Theorem~A supplies
strictness directly, without invoking innerness. We make no additional
finite-section claim for $\omega>\tfrac12$ beyond citing~\cite{Suzuki2012},
and the shift-radial Gram factor of Section~9.2 is the only factorization
theorem asserted in this paper. Extending the certificate of Theorem~B to
$\omega\in[\tfrac12,\omega_1]$ would give, via~\eqref{eq:energy-identity},
$D_{\omega,L}\succeq D_{1/2,L}\succ0$ and hence an independent
Paley--Wiener-free proof of strict contractivity on the first slab in that
range; floating diagnostics (Section~12.3) suggest this is available up to
$\omega_1\approx0.86$ in the even sector, but no such extension is claimed.

% =====================================================================
\section{Relation to the broader shifted-zeta program}
% =====================================================================

\subsection{String and screw-function realizations}

The shifted string Weyl function is
\begin{equation}
q_\omega(z)=\frac1{\sqrt{-z}}\frac{\xi'}{\xi}\Bigl(\tfrac12+\omega+\sqrt{-z}\Bigr),
\qquad\Re\sqrt{-z}>0 .
\end{equation}
Under the corresponding shifted zero-free hypothesis this is a Stieltjes
function and defines a positive Kre\u{\i}n string. Its spectral density is
\begin{equation}
\rho_\omega(u)=\frac1\pi\Re\frac{\xi'}{\xi}\Bigl(\tfrac12+\omega+iu\Bigr).
\end{equation}
The associated K\"all\'en--Lehmann superposition is reflection positive. This
is an exact inverse-spectral reformulation of the shifted zero-free condition
(compare the shifted screw functions $\Psi_\omega$
of~\cite[Sect.~11]{Suzuki2022screw}), but it is downstream of positivity. By
contrast, the first-slab Volterra system is specified directly from gamma and
prime coefficients and is positive without a zero-free assumption.

If $A_\omega(t)$ denotes Suzuki's integrated radial profile and
$\mathcal E_\omega=\int_0^\infty\abs{1-A_\omega(t)}^2\dd t$, then on the inner
side
\begin{equation}
\mathcal E_\omega=2\frac{\xi'}{\xi}\Bigl(\tfrac12+\omega\Bigr),
\qquad
\rho_\omega(0)=\frac{\mathcal E_\omega}{2\pi},
\qquad
\lim_{x\to\infty}\frac{m^{\mathrm{str}}_\omega(x)}{x}=\frac4{\mathcal E_\omega^2}.
\end{equation}
Thus the screw margin, phase density, boundary energy, and infrared string
density are different transforms of one scalar quantity. These companion
statements are developed in~\cite{Baker2026string,Baker2026margin,Baker2026depth}
and are not used here.

\subsection{Modular endpoint}

Let $\Lambda(s)=\pi^{-s/2}\Gamma(s/2)\zeta(s)$ and
$\varphi(s)=\Lambda(2s-1)/\Lambda(2s)$. At $\omega=\tfrac12$,
\begin{equation}
\Theta_{1/2}(u)=\frac{u-i}{u+i}\,\varphi\Bigl(\frac{1-iu}2\Bigr).
\end{equation}
The critical shift endpoint is therefore the modular cusp scattering
coefficient completed by one elementary Blaschke channel, in the standard
scattering-theoretic sense~\cite{GrahamZworski2003}. This supplies a geometric
anchor, but physical-line unitarity is unconditional; the RH-sensitive
content remains off-line causal analyticity and contractivity.

\subsection{Radial reflection positivity and holographic language}

The term \emph{radial reflection positivity} is literal here: one logarithmic
coordinate is reflected, and the boundary form is the contraction defect of a
causal radial propagator. The identity
$D_{\omega,L}=2\int_0^\omega V_{s,L}^{*}Q_{s,L}V_{s,L}\dd s$ is a positive-bulk
reconstruction of the boundary defect. The word \emph{holographic} is useful
only in this transfer-system sense. The construction is not presently a local
Poincar\'e--Einstein geometry, and the boundary object is not generally a
homogeneous CFT two-point function. A closer PDE precedent is Burnol's causal
Klein--Gordon realization of the classical Hankel transform~\cite{Burnol2005}.

\subsection{Where arithmetic first enters}

For $\log2<L<\log3$ the kernel becomes
\begin{equation}
k_\omega(t)=\kappa_\omega(t)+\frac{c_\omega(2)}{\sqrt2}\kappa_\omega(t-\log2).
\end{equation}
The new term is a delayed copy of the archimedean impulse. Extending the
positive energy identity across this update is the first genuinely arithmetic
problem. The present theorem supplies the base storage operator against which
that delayed layer must be tested.

% =====================================================================
\section{Alternative checks and superseded proof routes}
% =====================================================================

\subsection{Compact Mosco limit}

The normalized defect channels $\omega^{-1}\langle f,(I\mp H_{\omega,L})f\rangle$
converge as $\omega\downarrow0$ to the even and odd restrictions of $Q_{0,L}$.
Uniform logarithmic coercivity prevents high-frequency escape, and compact
Mosco convergence combines with the central strict gaps to give a qualitative
small-shift theorem. This route remains an independent analytic explanation
of the singular endpoint. It is no longer needed for the complete interval
theorem, because the shift-energy proof uses the strong multiplier limit
$V_s\to I$ directly.

\subsection{Direct transfer certificates}

Earlier calculations certified the defect channels directly through a
parity-scaled matrix-valued Taylor model and centered endpoint charts. These
computations supplied substantial numerical and interval evidence, but the
finite-$\omega$ defect matrices mix parity and require more elaborate
conditioning than the generator form. The generator certificate is preferable
because the generator remains parity diagonal; all shift dependence is a
compact kernel entire in $\omega^2$; the degree-four remainder is below
$1.54\times10^{-12}$; the full interval, including both endpoints, is covered
uniformly; and the energy identity converts generator positivity into the
desired defect inequality. The earlier certificates are retained in the
research archive as independent cross-checks, but not as part of the main
proof.

% =====================================================================
\section{Limitations and next problems}
% =====================================================================

\subsection{A fully dual-library certificate}

The clean-room audit independently reconstructs the mathematics and interval
$LDL^{\mathsf T}$ calculation. Its central special-function path uses
arbitrary precision with a conservative audit radius rather than a second
formal interval library. A publication repository would be stronger if the
central digamma/Bessel quadrature were also implemented in a second validated
interval environment. The current agreement makes this an audit enhancement,
not a practical blocker.

\subsection{The second slab}

The immediate mathematical target is the delayed $n=2$ layer. Natural
approaches include differentiating the complete delayed propagator, deriving a
storage update or Redheffer-star-product inequality, and testing whether the
delayed layer preserves passivity.

The archimedean continuation cannot be assumed positive throughout this next
slab. A 24-mode floating diagnostic, independently rerun but not promoted to
an interval certificate, places the first odd and even crossings of the
prime-free generator near
\begin{equation}
L^{*}_{\mathrm{odd}}\approx0.7429,\qquad L^{*}_{\mathrm{even}}\approx0.8199 .
\end{equation}
In particular, the odd sector is already negative at $L=0.75$, only about
seven percent beyond $\log2$, and both sectors are negative before $\log3$.
These values concern the prime-free archimedean continuation; they do not
contradict the prime-active full-Weil certificates of~\cite{Zhu2026}. They
change the logic of the second-slab project: the delayed $n=2$ term cannot be
treated over the whole slab merely as a small correction to a positive base.
It must be analyzed as a possible source of the missing positivity, or as the
source of a precise obstruction.

At the fixed first-slab horizon $L=\log2$, the same floating diagnostic places
the first even generator crossing near $\omega=0.8611$ and the odd crossing
near $\omega=1.8047$. Those estimates suggest a possible extension of the
generator certificate beyond the critical interval (Section~9.3), but such an
extension is outside Theorem~B and should be stated only after a separate
interval run.

The next steps are therefore: derive the exact differentiated $n=2$ delayed
generator; isolate its signed contribution relative to the archimedean
continuation; seek a storage or Redheffer inequality that allows the delayed
layer to restore, rather than merely preserve, positivity; certify
representative horizon slices in interval arithmetic; and identify the first
rigorous obstruction if the repair mechanism fails.

\subsection{Euler-product network and geometry}

The coefficient train satisfies
\begin{equation}
\sum_{n\ge1}\frac{c_\omega(n)}{\sqrt n}n^{-s}
=\frac{\zeta(s+\tfrac12-\omega)}{\zeta(s+\tfrac12+\omega)} .
\end{equation}
A passive-network composition law reflecting this Euler product would be a
natural global continuation of the first-slab theorem. Separately, one may
seek a local PDE or boundary-triple realization of the completed-gamma
quotient whose energy estimate produces $Q_{s,L}$ directly.

% =====================================================================
\section{Verification status}\label{sec:verification}
% =====================================================================

This section replaces the private author-verification notice of the working
drafts. It records what has been checked, by what means, and what has not.
The manuscript has been through two adversarial referee-style review rounds
with independent numerical recomputation of every reported constant, and one
clean-room reimplementation of the certificate (Section~\ref{sec:audit}). No
error survives in the main proof chain. What the manuscript has \emph{not}
had is line-by-line derivation of the central analytic steps by a human
reader who will stand behind them; the deposit of this preprint is not a
claim that this has happened.

\subsection*{Independently reproduced from the printed definitions}

Two further recomputations, independent of both the primary Arb code and the
clean-room audit---one made during the referee-style review rounds, one made
while preparing this preprint---reproduced the following from the definitions
printed here:
all algebraic identities in Sections~2.2, 3.1, 3.3, 4.1, 6.1, 6.2, 10.2 and
10.4; the Laplace transform of the explicit impulse~\eqref{eq:kappa-explicit}
against~\eqref{eq:K} (to $1.5\times10^{-17}$ at $\omega=\tfrac12$ and
$8\times10^{-11}$ at $\omega=0.3$, the latter limited by quadrature at the
$t^{\omega-1}$ singularity); the multiplier identity of
Proposition~\ref{prop:bounded-difference} to 30 digits at $\omega=0.1$, $0.3$
and $\tfrac12$; each step of its proof (digamma integral representation,
cancellation of the divergent summands, the rational kernel
$4\cosh(t/2)(\cosh\omega t-1)$, and the bridge identity for $R$); the
endpoint Lemma~\ref{lem:endpoint}, whose corrected formula agrees with
$Q_0+C_{1/2}$ for test functions to the quadrature floor
($\sim10^{-8}$) while the uncorrected pointwise-limit formula is off by
$0.2402$; the complete chain of Section~\ref{sec:kernel-bound} ($t_0$, the
closed form of $h$, the sign analysis of $N$, $t_m$, $h(t_m)$,
$h(\log2-t_0)$, $\int_0^{\log2}\abs h=0.012293050401142319559\ldots$, and the
inequality $L/2>t_0$); the tail bound~\eqref{eq:dtail} and both Schur losses
to 17 digits; the Taylor remainder $1.532967865231399\times10^{-12}$; the
head minima and margins of Section~7.2 by an independent sixteen-mode
Legendre reconstruction, consistent with the reported values to within the
expected interval-widening penalty ($2.26\times10^{-5}$ even,
$8.65\times10^{-6}$ odd); the coercivity constants of
Remark~\ref{rem:not-m}; and all reference metadata.

\subsection*{Checked against Suzuki's statements}

The transfer normalization~\eqref{eq:K} and the layer
decomposition~\eqref{eq:layer} were checked against the displayed statements
of~\cite{Suzuki2012} as follows: Suzuki's factor $\gamma(s)$ is
$\tfrac12s(s-1)\pi^{-s/2}\Gamma(s/2)=\Ginf(s)$; his Mellin
formula~\cite[Sect.~3.1]{Suzuki2012} for $g_\omega$ is~\eqref{eq:K} after
$s=p+\tfrac12$; \eqref{eq:layer} is~\cite[eq.~(2.3)]{Suzuki2012} in
logarithmic coordinates; and~\cite[Prop.~2.1]{Suzuki2012} is consistent with
these two facts and the Dirichlet series of $c_\omega$, with matching domains
of convergence. The identification of $\Hsf_{\omega,a}$ with the centered
interval of length $2\log a$ (Section~2.2) is derived from
\cite[eqs.~(2.5)--(2.6)]{Suzuki2012}. These comparisons were made against
automated readings of the published text and should be confirmed by a human
reader against the paper itself; that confirmation is the single most
important outstanding check, since a factor or shift there would move the
$\log2$ threshold.

\subsection*{Not yet human-verified}

\begin{enumerate}[nosep,label=(\alph*)]
\item \emph{Elementary derivations on paper}: Proposition~\ref{prop:bounded-difference}
  including its endpoint extension; Lemma~\ref{lem:endpoint};
  Lemma~\ref{lem:row-max}; the Volterra identities~\eqref{eq:volterra-identities}.
\item \emph{Functional analysis of Section~8}: the operator-norm $C^1$
  dependence of $s\mapsto V_{s,L}$ (the multiplier bounds of Section~8.1 are
  stated with $O(\cdot)$ and no constants, and are the weakest-supported
  claims in the manuscript); the weighted-to-unweighted transfer; the
  closure step~\eqref{eq:Q-pairing}; the improper strong-operator integral
  of Section~8.5. Failure of these would leave Theorem~C unproved and
  Theorem~A resting on the routes of Section~11, not make Theorem~A false.
\item \emph{Analytic number theory}: human confirmation of the normalization
  chain against~\cite{Suzuki2012} (above); an originality search for the
  off-center kernel identity of Proposition~\ref{prop:bounded-difference} and
  the shift-radial energy identity against~\cite{ConnesConsani2020,Zhu2026,Kim2026,Groskin2026}
  and newer work in this active neighborhood.
\item \emph{Repository}: the inherited central certificate (the Binet bound
  $w_0(\tau)\ge\beta_*$ for $\abs\tau\ge30$ and the spherical-Bessel
  Legendre-tail majorants behind~\eqref{eq:central-couplings}), which both the
  primary and audit implementations take as given; the interval
  $LDL^{\mathsf T}$ and box-construction code; a full primary-versus-audit row
  diff rather than a comparison of summary minima; and a version-pinned rerun
  of both implementations in clean environments with archived logs.
\end{enumerate}
Optional strengthenings (not verification): promote the positivity-horizon
crossings of Section~12.3 to interval certificates; extend the box sweep to
$\omega\approx0.85$ (Section~9.3); double the box count to $512$, since the
box radius ($2.42\times10^{-5}$) is now the second-largest erosion of the even
margin after the Schur loss.

% =====================================================================
\appendix
\section{Normalization ledger}
% =====================================================================

The manuscript uses
\begin{gather}
\xi(s)=\tfrac12s(s-1)\pi^{-s/2}\Gamma(s/2)\zeta(s),\qquad
\Ginf(s)=\tfrac12s(s-1)\pi^{-s/2}\Gamma(s/2),\\
\Theta_\omega(z)=\frac{\xi(\tfrac12-\omega-iz)}{\xi(\tfrac12+\omega-iz)},\qquad
K_\omega(p)=\frac{\Ginf(\tfrac12+p-\omega)}{\Ginf(\tfrac12+p+\omega)},\\
Q_{\omega,L}=\tfrac12\bigl(A_{\omega,L}+A_{\omega,L}^{*}\bigr),\qquad
L_{\max}=\log2,\qquad T_0=\tfrac12\log2 .
\end{gather}
Suzuki's multiplicative compression $\Hsf_{\omega,a}=P_a\Hsf_\omega P_a$
corresponds to $L=2\log a$; the compact-window convention of~\cite{Zhu2026}
corresponds to $L=2\LZ$. For the shifted string,
\begin{equation}
\rho_\omega(u)=\frac1\pi\Re\frac{\xi'}{\xi}(\tfrac12+\omega+iu),\qquad
q_\omega(z)=\frac1{\sqrt{-z}}\frac{\xi'}{\xi}(\tfrac12+\omega+\sqrt{-z}).
\end{equation}
The factor $Q=B/2$, where $B=A+A^{*}$, is essential when comparing the central
Weil certificate, perturbation matrices, and shift-energy identity.

% =====================================================================
\section{Reproducibility}
% =====================================================================

The certificate scripts described below are not part of this deposit; they
are available from the author on request, and a code release with pinned
environments and archived logs is planned once the repository items of
Section~\ref{sec:verification}(d) are complete.

\subsection*{B.1 Primary certificate}

The primary generator certificate is produced by the script
\begin{quote}
\texttt{investigation12\_offcenter\_weil\_generator.py},
\end{quote}
with the inherited central Weil implementation in
\begin{quote}
\texttt{project\_support/investigation7\_first\_slab\_certificate.py}:
\end{quote}
\begin{verbatim}
PYTHONPATH=project_support:. \
python3 investigation12_offcenter_weil_generator.py \
  --certify --boxes 256 --output INVESTIGATION_12_results.csv
python3 verify_investigation12_results.py
\end{verbatim}
Expected summary:
\begin{verbatim}
certified boxes: 512
kernel norm < 0.012294
Taylor remainder < 1.54e-12
minimum interval LDL pivot > 0.0232681
minimum certified Weyl margin > 0.00061163
\end{verbatim}
The 724 result rows consist of 6 Laplace-transform identity checks, 1 kernel
norm row, 1 Taylor-remainder row, 202 floating generator samples (101 shifts
in two parity sectors), 2 relative-endpoint checks, and 512 certified
parameter boxes.

\subsection*{B.2 Independent audit}

The audit is produced by the script
\begin{quote}
\texttt{investigation14\_independent\_generator\_audit.py}:
\end{quote}
\begin{verbatim}
python3 investigation14_independent_generator_audit.py \
  --primary INVESTIGATION_12_results.csv \
  --output INVESTIGATION_14_results.csv \
  --matrices INVESTIGATION_14_matrices.csv
python3 verify_investigation14_results.py
\end{verbatim}
Expected summary:
\begin{verbatim}
audit rows: 783
matrix rows: 1360
maximum central quadrature discrepancy: 1.693e-15
maximum primary floating discrepancy: 4.375e-13
maximum direct-kernel discrepancy: 1.593e-13
minimum directed-interval margin: 6.116514636999e-04
minimum directed-interval pivot: 2.326829938890e-02
\end{verbatim}
The 783 audit rows split as follows: 2 central-quadrature convergence rows,
1 kernel norm, 1 Taylor remainder, 24 Legendre-correlation comparisons,
5 profile-series comparisons, 36 direct-kernel comparisons, 202 primary
floating comparisons, and 512 directed-interval boxes. The 1360 matrix rows
are the lower triangles of ten $16\times16$ coefficient matrices: five Taylor
matrices in each of two parity sectors, with 136 stored entries per matrix.

\subsection*{B.3 Analytic diagnostics}

The shift-energy transfer diagnostics are generated by
\begin{quote}
\texttt{investigation13\_operator\_diagnostics.py}
\end{quote}
and checked by \texttt{verify\_investigation13\_results.py}. They verify the transfer
derivative, pointwise multiplier limit, high-frequency regularization, and
local kernel tail on a high-precision sample grid. These checks are not
ingredients of the analytic proof.

% =====================================================================
\section{Proof-status ledger}
% =====================================================================

\begin{center}
\small
\begin{tabular}{@{}p{0.55\textwidth}p{0.40\textwidth}@{}}
\toprule
Statement & Status\\
\midrule
First-slab Volterra reduction & exact\\
Completed-gamma transfer & exact; matches~\cite[Sect.~3.1]{Suzuki2012}\\
Sonine convolution identity & exact\\
Formula~\eqref{eq:a-mult} for $a_\omega$ & exact\\
Fourier-plus-residue endpoint term & exact distributional correction\\
Bounded-kernel decomposition $Q_\omega=Q_0+C_\omega$ & exact\\
Common logarithmic form domain & exact by bounded form perturbation\\
Central prime-free Weil positivity & Arb-certified\\
Sharp endpoint-row Schur bound & analytic monotonicity plus Arb scalar enclosure\\
Degree-four shift expansion and remainder & analytic plus Arb enclosure\\
All 512 generator boxes & Arb-certified\\
Infinite Legendre-tail closure & analytic bounds plus certified head\\
Block-Schur coercivity on $\Dlog$ & exact (Lemma~\ref{lem:block-schur})\\
Independent implementation audit & passed\\
Strong limit $V_s\to I$ & analytic Hardy-multiplier proof\\
Closed-form shift-energy identity & analytic proof\\
Strict contractivity & analytic consequence of generator positivity and compactness\\
Shift-radial Gram factor & exact consequence\\
Strict contractivity for $\omega>\tfrac12$ & Suzuki~\cite[Lemma~4.4]{Suzuki2012}, not reproved here\\
Positivity beyond the first slab & open\\
New zero-free region or RH & not claimed\\
\bottomrule
\end{tabular}
\end{center}

% =====================================================================
\begin{thebibliography}{99}
\small

\bibitem{Suzuki2012}
M.~Suzuki,
\emph{A canonical system of differential equations arising from the Riemann
zeta-function}, arXiv:1204.1827.

\bibitem{Suzuki2022screw}
M.~Suzuki,
\emph{Aspects of the screw function corresponding to the Riemann
zeta-function}, arXiv:2206.03682.

\bibitem{Suzuki2026screw}
M.~Suzuki,
\emph{Weil's quadratic form via the screw function}, arXiv:2606.09096.

\bibitem{ConnesConsani2020}
A.~Connes and C.~Consani,
\emph{Weil positivity and trace formula, the archimedean place},
Selecta Math. (N.S.) 27 (2021), no.~4, Paper No.~77; arXiv:2006.13771.

\bibitem{Zhu2026}
X.~Zhu,
\emph{Weil positivity in compact windows: a finite reduction, certified
two-sided bounds, and a Landau--Widom decay law}, arXiv:2608.24827 (v2).

\bibitem{Burnol2005}
J.-F.~Burnol,
\emph{Spacetime causality in the study of the Hankel transform},
arXiv:math/0509619.

\bibitem{Lagarias1999}
J.~C.~Lagarias,
\emph{On a positivity property of the Riemann $\xi$-function},
Acta Arith. 89 (1999), 217--234; correction, Acta Arith. 116 (2005), 293--294.

\bibitem{Luchko2021}
Y.~Luchko,
\emph{Operational calculus for the general fractional derivatives with the
Sonine kernels}, arXiv:2103.00549.

\bibitem{OS1973}
K.~Osterwalder and R.~Schrader,
\emph{Axioms for Euclidean Green's functions},
Comm. Math. Phys. 31 (1973), 83--112.

\bibitem{GrahamZworski2003}
C.~R.~Graham and M.~Zworski,
\emph{Scattering matrix in conformal geometry},
Invent. Math. 152 (2003), 89--118; arXiv:math/0109089.

\bibitem{DLMF}
NIST Digital Library of Mathematical Functions, Chapter~5: Gamma Function,
Sections~5.9, 5.11 and 5.15, \url{https://dlmf.nist.gov/5}.

\bibitem{Kim2026}
T.~Kim, Y.~Hong, M.~Kim, S.~Choi, J.~Jang, and M.~Kim,
\emph{A numerical realization of Suzuki's Weil-quadratic-form operator: the
archimedean spectral law, its universality, and an operator form of Weil's
positivity criterion}, arXiv:2607.24830.

\bibitem{Groskin2026}
A.~Groskin,
\emph{A finite Guinand--Weil dictionary and archimedean tail order for the
truncated Weil quadratic form}, arXiv:2607.02828.

\bibitem{Baker2026string}
E.~Baker,
\emph{The shifted zeta string: an unconditional inverse-spectral family and
its RH endpoint}, manuscript, August 2026.

\bibitem{Baker2026margin}
E.~Baker,
\emph{The margin in the shifted screw-function criterion for zero-free
half-planes}, manuscript, August 2026.

\bibitem{Baker2026depth}
E.~Baker,
\emph{Detection depth of spectral defects in shifted zeta strings},
manuscript, August 2026.

\end{thebibliography}

\end{document}
